\section{总结与延伸}

\subsection{核心结论}

\begin{enumerate}
\item Manus 的故事要从 App Store 早期个人开发者经历读起，而不是只从 Meta 收购读起。
\item Peak 的学习路径由具体产品需求牵引：浏览器预加载引向 NLP，NLP 引向知识结构化和 Agent。
\item Manus 从 0 到 1 亿 ARR 的关键，是从 demo 进入真实工作流和收入飞轮。
\item AI 更像制造业，意味着 Agent 公司要面对质量、成本、产能、生态配套和交付。
\item Agent 继续爆发需要身份、支付、权限、反爬和审计等基础设施配合。
\item Manus 最大风险不是单纯大厂竞争，而是失去特色和变得复杂。
\item Agent 不一定竞争用户时长，它可能竞争后台执行、授权、信任和生产力结果。
\item 新加坡和全球化选择，本质上是客户、合规、支付、数据和组织协作问题。
\item 保持简单不是少做，而是让新增功能增强核心价值。
\end{enumerate}

\subsection{开放问题}

最后保留开放问题，是因为 Manus 已经进入新的公司阶段，但 Agent 产品和市场形态还远没有收敛。

\begin{itemize}
\item Meta 收购后，Manus 的简单性和特色能否保持？
\item C 端 Agent 是否会真正越过生产力工具门槛，成为普通白领的主力应用？
\item 支付、权限、人机验证和网站开放，会不会成为 Agent 爆发的关键基础设施？
\item Agent 产品的稳态，会不会不同于移动互联网的用户时长竞争？
\item “AI 像制造业”是否意味着未来 Agent 公司会更像供应链/工厂，而不是传统软件公司？
\end{itemize}

\subsection{拓展阅读}

\begin{itemize}
\item EP095 Manus 创始人肖弘访谈：理解 Manus 的前史和博弈观。
\item EP130 张月光访谈：AI Native 产品、Agent 和 AI 朋友。
\item EP139 Agent 技术史：Agent 的技术谱系和产品化漏斗。
\item Open IE、Word2Vec、知识图谱相关材料：理解 Peak 早期技术路径。
\item Coding Agent、Sierra、Replit、Lovable、Claude Code 等案例：理解 Agent 市场分化。
\end{itemize}

\begin{importantbox}{最后的判断}
Manus 的最后访谈最值得留下的，不是“被谁收购”这个结局，而是一个 Agent 产品团队在增长前夜的自我提醒：如果增长让产品失去特色，成功本身就会变成新的风险。
\end{importantbox}

\end{document}
